\SecondArticle{9}
Den Beweis des ersten Satzes werden wir so abschließen, dass wir zeigen: Haben $Y$ und $Y'$ keinen gemeinsamen Teiler, so kann gewiss nicht $P=0$ gelten. Dazu nehmen wir zunächst an, nach den Vorschriften von Artikel 2 seien zwei Polynome in der Unbestimmten $x$, nämlich $fx$ und $\varphi x$, gefunden worden, sodass die identische Gleichung
\[
fx\, .\,Y+\varphi x\, .\,Y'=1
\]
gilt. Wir schreiben sie hier in der Form
\[
fx\, .\,\upsilon+\varphi x\, .\,\upsilon'=1+fx\, .\,(\upsilon-Y)+\varphi x\, .\,\frac{\mathrm d(\upsilon-Y)}{\mathrm dx}
\]
oder, da
\[
\begin{aligned}
\upsilon'={}&(x-b)(x-c)(x-d)\, .\, .\, .\, .\\
&+(x-a)\, .\,\frac{\mathrm d[(x-b)(x-c)(x-d)\, .\, .\, .\, .]}{\mathrm dx}
\end{aligned}
\]
gilt, in der folgenden Form:
\[
\begin{aligned}
&\varphi x\, .\,(x-b)(x-c)(x-d)\, .\, .\, .\, .\\
&+\varphi x\, .\,(x-a)\, .\,\frac{\mathrm d[(x-b)(x-c)(x-d)\, .\, .\, .\, .]}{\mathrm dx}\\
&+fx\, .\,(x-a)(x-b)(x-c)(x-d)\, .\, .\, .\, .\\&=1+fx\, .\,(\upsilon-Y)+\varphi x\, .\,\frac{\mathrm d(\upsilon-Y)}{\mathrm dx}
\end{aligned}
\]
\Needspace{8\baselineskip}Der Kürze halber bezeichnen wir das Polynom
\[
fx\, .\,(y-Y)+\varphi x\, .\,\frac{\mathrm d(y-Y)}{\mathrm dx}
\]
\SourceLeaf{050}
in den Unbestimmten $x,l',l'',l'''$ usw.
\[
\text{mit}\quad F(x,l',l'',l'''\text{ usw.})
\]
Damit gilt identisch
\[
\begin{gathered}1+fx\, .\,(\upsilon-Y)+\varphi x\, .\,\frac{\mathrm d(\upsilon-Y)}{\mathrm dx}\\=1+F(x,\lambda',\lambda'',\lambda'''\text{ usw.})\end{gathered}
\]
Wir erhalten also die identischen Gleichungen [1]
\[
\begin{aligned}
&\varphi a\, .\,(a-b)(a-c)(a-d)\, .\, .\, .\, .\\&=1+F(a,\lambda',\lambda'',\lambda'''\text{ usw.})\\
&\varphi b\, .\,(b-a)(b-c)(b-d)\, .\, .\, .\, .\\&=1+F(b,\lambda',\lambda'',\lambda'''\text{ usw.})\\
&\varphi c\, .\,(c-a)(c-b)(c-d)\, .\, .\, .\, .\\&=1+F(c,\lambda',\lambda'',\lambda'''\text{ usw.})\\
&\text{usw.}
\end{aligned}
\]
Nehmen wir nun an, das Produkt aus allen Ausdrücken
\[
\begin{gathered}
1+F(a,l',l'',l'''\text{ usw.})\\
1+F(b,l',l'',l'''\text{ usw.})\\
1+F(c,l',l'',l'''\text{ usw.})\\
\text{usw.}
\end{gathered}
\]
das ein Polynom in den Unbestimmten $a,b,c$ usw., $l',l'',l'''$ usw. und hinsichtlich $a,b,c$ usw. symmetrisch ist, werde durch
\[
\psi(\lambda',\lambda'',\lambda'''\text{ usw.},\ l',l'',l'''\text{ usw.})
\]
dargestellt. Dann ergibt die Multiplikation sämtlicher Gleichungen [1] die neue identische Gleichung [2]
\[
\pi\varphi a\, .\,\varphi b\, .\,\varphi c\, .\, .\, .\, .=\psi(\lambda',\lambda'',\lambda'''\text{ usw.},\ \lambda',\lambda'',\lambda'''\text{ usw.})
\]
Da ferner das Produkt $\varphi a\, .\,\varphi b\, .\,\varphi c\, .\, .\, .\, .$ die Unbestimmten $a,b,c$ usw. symmetrisch enthält, lässt sich offenbar ein Polynom in den Unbestimmten $l',l'',l'''$ usw. finden, das durch die Substitutionen $l'=\lambda',l''=\lambda'',l'''=\lambda'''$ usw. in $\varphi a\, .\,\varphi b\, .\,\varphi c\, .\, .\, .\, .$ übergeht. Dieses Polynom sei $t$. Dann gilt ebenfalls \mbox{identisch [3]}
\[
pt=\psi(l',l'',l'''\text{ usw.},\ l',l'',l''')
\]
denn die Substitutionen $l'=\lambda',l''=\lambda'',l'''=\lambda'''$ usw. führen diese Gleichung in die identische Gleichung [2] über.

Aus der Definition der Funktion $F$ selbst folgt nun identisch
\SourceLeaf{051}
\[
F(x,L',L'',L'''\text{ usw.})=0
\]
Daher gilt auch identisch
\[
\begin{gathered}
1+F(a,L',L'',L'''\text{ usw.})=1\\
1+F(b,L',L'',L'''\text{ usw.})=1\\
1+F(c,L',L'',L'''\text{ usw.})=1\\
\text{usw.}
\end{gathered}
\]
und somit ebenfalls identisch
\[
\psi(\lambda',\lambda'',\lambda'''\text{ usw.},\ L',L'',L'''\text{ usw.})=1
\]
also auch identisch [4]
\[
\psi(l',l'',l'''\text{ usw.},\ L'.\,L'',L'''\text{ usw.})=1
\]
Verbinden wir die Gleichungen [3] und [4] und setzen $l'=L'$, $l''=L'',l'''=L'''$ usw., so erhalten wir daher [5]
\[
PT=1
\]
wobei $T$ den Wert des Polynoms $t$ bei diesen Substitutionen bezeichnet. Da dieser Wert notwendig eine endliche Größe ist, kann $P$ gewiss nicht $=0$ sein. Q. E. D.

\Needspace{18\baselineskip}\SecondArticle{10}
Aus dem Vorhergehenden ist nun klar, dass sich jedes Polynom $Y$ in einer Unbestimmten $x$, dessen Determinante $=0$ ist, in Faktoren zerlegen lässt, von denen keiner die Determinante 0 hat. Bestimmt man nämlich den größten gemeinsamen Teiler der Polynome $Y$ und $\frac{\mathrm dY}{\mathrm dx}$, so ist das erste damit in zwei Faktoren zerlegt. Hat einer dieser Faktoren\HistoricalNote{11} wiederum die Determinante 0, so wird er auf dieselbe Weise in zwei Faktoren zerlegt. So fahren wir fort, bis $Y$ schließlich in Faktoren zerlegt ist, von denen keiner die Determinante 0 hat.

Ferner sieht man leicht, dass sich unter diesen Faktoren, in die $Y$ zerlegt wird, wenigstens
\SourceLeaf{052}
einer befinden muss, der folgende Eigenschaft hat: Unter den Faktoren der Zahl, die seinen Grad angibt, kommt die Zahl zwei nicht öfter vor als unter den Faktoren der Zahl $m$, die den Grad des Polynoms $Y$ angibt. Setzt man also $m=k\, .\,2^\mu$, wobei $k$ eine ungerade Zahl bezeichnet, so findet sich unter den Faktoren des Polynoms $Y$ wenigstens einer vom Grad $k'\, .\,2^\nu$, wobei auch $k'$ ungerade ist und entweder $\nu=\mu$ oder $\nu<\mu$ gilt. Die Richtigkeit dieser Behauptung folgt von selbst daraus, dass $m$ die Summe der Zahlen ist, welche die Grade der einzelnen Faktoren von $Y$ angeben.

\SecondArticle{11}
Bevor wir weitergehen, erläutern wir einen Ausdruck, dessen Einführung bei allen Untersuchungen über symmetrische Funktionen sehr nützlich ist und der auch uns sehr gelegen kommt. Es sei $M$ eine Funktion einiger der Unbestimmten $a,b,c$ usw.; $\mu$ sei dabei die Anzahl derjenigen, die in den Ausdruck $M$ eingehen, ohne Rücksicht auf andere Unbestimmte, die $M$ etwa noch enthält. Vertauscht man diese $\mu$ Unbestimmten auf alle möglichen Arten sowohl untereinander als auch mit den übrigen $m-\mu$ aus $a,b,c$ usw., so entstehen aus $M$ andere, $M$ ähnliche Ausdrücke. Insgesamt erhält man
\[
m(m-1)(m-2)(m-3)\, .\, .\, .\, .\, .\,(m-\mu+1)
\]
Ausdrücke einschließlich $M$ selbst. Ihre Gesamtheit nennen wir kurz die \textit{Gesamtheit aller $M$.} Damit ist von selbst klar, was die Summe aller $M$, das Produkt aus allen $M$ usw. bedeutet. So heißt zum Beispiel $\pi$ das Produkt aus allen $a-b$, $\upsilon$ das Produkt aus allen $x-a$ und $\upsilon'$ die Summe aller $\frac{\upsilon}{x-a}$ usw.

Ist $M$ hinsichtlich einiger der $\mu$ darin enthaltenen Unbestimmten symmetrisch, so ändern deren Vertauschungen untereinander die Funktion $M$ nicht. In der Gesamtheit aller $M$ kommt daher jedes Glied mehrfach vor, und zwar $1\, .\,2\, .\,3\, .\, .\, .\, .\, .\,\nu$ Mal, wenn $\nu$ die Anzahl der Unbestimmten ist, hinsichtlich deren $M$ symmetrisch ist. Ist aber $M$ nicht nur hinsichtlich $\nu$ Unbestimmter symmetrisch, sondern außerdem hinsichtlich $\nu'$ anderer und ferner hinsichtlich $\nu''$ anderer usw., so bleibt $M$ unverändert, ob nun zwei der ersten $\nu$ Unbestimmten untereinander vertauscht werden oder zwei der zweiten $\nu'$ oder zwei der dritten $\nu''$ usw. Es entsprechen somit stets
\[
1\, .\,2\, .\,3\, .\, .\, .\, .\, .\,\nu\, .\quad1\, .\,2\, .\,3\, .\, .\, .\, .\, .\,\nu'\, .\quad1\, .\,2\, .\,3\, .\, .\, .\, .\, .\,\nu''\text{ usw.}
\]
\SourceLeaf{053}
Vertauschungen identischen Gliedern. Behalten wir von diesen identischen Gliedern jeweils nur eines, so erhalten wir insgesamt
\[
\frac{m(m-1)(m-2)(m-3)\, .\, .\, .\, .\, .\,(m-\mu+1)}{1\, .\,2\, .\,3\, .\, .\, .\, .\, .\,\nu\, .\,1\, .\,2\, .\,3\, .\, .\, .\, .\, .\,\nu'\, .\,1\, .\,2\, .\,3\, .\, .\, .\, .\, .\,\nu''\text{ usw.}}
\]
Glieder. Ihre Gesamtheit nennen wir die \textit{Gesamtheit aller $M$ mit ausgeschlossenen Wiederholungen}, um sie von der \textit{Gesamtheit aller $M$ mit zugelassenen Wiederholungen} zu unterscheiden. Wird nichts ausdrücklich bemerkt, so verstehen wir stets, dass die Wiederholungen zugelassen sind.

Ferner sieht man leicht, dass die Summe aller $M$ oder das Produkt aus allen $M$ oder allgemein irgendeine symmetrische Funktion aller $M$ stets eine symmetrische Funktion der Unbestimmten $a,b,c$ usw. ist, ob Wiederholungen zugelassen oder ausgeschlossen werden.

\SecondArticle{12}
Nun betrachten wir, mit $u,x$ als Unbestimmten, das Produkt aus allen $u-(a+b)x+ab$ mit ausgeschlossenen Wiederholungen und bezeichnen es mit $\zeta$. Somit ist $\zeta$ das Produkt aus den folgenden \mbox{$\frac12m(m-1)$} Faktoren:
\[
\begin{gathered}
u-(a+b)x+ab\\
u-(a+c)x+ac\\
u-(a+d)x+ad\\
\text{usw.}\\
u-(b+c)x+bc\\
u-(b+d)x+bd\\
\text{usw.}\\
u-(c+d)x+cd\\
\text{usw. usw.}
\end{gathered}
\]
Da diese Funktion die Unbestimmten $a,b,c$ usw. symmetrisch enthält, lässt sich ein Polynom in den Unbestimmten $u,x,l',l'',l'''$ usw. angeben, das wir mit $z$ bezeichnen und das in $\zeta$ übergeht, wenn anstelle der Unbestimmten $l',l'',l'''$ usw. $\lambda',\lambda'',\lambda'''$ usw. eingesetzt werden. Schließlich bezeichnen wir mit $Z$ die Funktion allein der Unbestimmten $u,x$, in die $z$ übergeht, wenn den Unbestimmten $l',l'',l'''$ usw. die bestimmten Werte $L',L'',L'''$ usw. gegeben werden.

Diese drei Funktionen $\zeta,z,Z$ lassen sich als Polynome vom Grad\\
\mbox{$\frac12m(m-1)$} in der Unbestimmten $u$ mit unbestimmten Koeffizienten auffassen. Diese
\SourceLeaf{054}
Koeffizienten sind
\[
\begin{aligned}
&\text{für }\zeta,\quad\text{Funktionen der Unbestimmten }x,a,b,c\text{ usw.}\\
&\text{für }z,\quad\text{Funktionen der Unbestimmten }x,l',l'',l'''\text{ usw.}\\
&\text{für }Z,\quad\text{Funktionen allein der Unbestimmten }x.
\end{aligned}
\]
\Needspace{12\baselineskip}Die einzelnen Koeffizienten von $z$ gehen durch die Substitutionen $l'=\lambda',l''=\lambda'',l'''=\lambda'''$ usw. in die Koeffizienten von $\zeta$ und durch die Substitutionen $l'=L',l''=L'',l'''=L'''$ usw. in die Koeffizienten von $Z$ über. Was wir eben über die Koeffizienten gesagt haben, gilt ebenso für die Determinanten der Funktionen $\zeta,z,Z$. Diese wollen wir nun näher untersuchen, und zwar mit dem Ziel, den folgenden Satz zu beweisen.

\textsc{Satz.} \textit{Gilt nicht $P=0$, so kann die Determinante der Funktion $Z$ gewiss nicht identisch $=0$ sein.}

